{\Large \bf Section 6: The Two-Body Problem}

\vskip 3mm
We're going to end this series of  lectures by looking at something deceptively simple: the motion of two particles, interacting through some potential. 

\para
Now that seems rather limited. Surely it would be more exciting to look at many particles interacting with other. And it is. But that also turns out to be hard -- hard as in no-one-knows-how-to-do-it-hard. Looking at three or more particles leads us into the world of chaos theory and numerical simulation. By contrast, understanding the motion of two particles -- the so-called two-body problem -- is something that we can, and will, solve on the board.

\para
It's also important. Understanding the two-body problem historically led to two of the biggest breakthroughs in the history of science. It was how we understood gravity, with Newton postulating his inverse square law to agree with the observations of the astronomer Kepler. And it's how we understood the structure of atoms, with Rutherford understanding how his own experiment could be interpreted as the atom having a tiny nucleus, surrounded by electrons. We're going to understand both these topics in these lectures. 

\section{The Two-Body Problem}

{\co We will consider two particles, interacting through a potential  $V_{12}=V_{21}=V$, such that
%
\be V = V(|\bx_1-\bx_2|)\ .\label{2bodypot}\ee
%
}
We saw earlier in these lectures that ensures that Newton's third law is obeyed. 

\para
We will first derive a number of results about the general two-body problem but, ultimately, we will focus on a very special force with potential $V\sim 1/r$ with $r$ the separation between particles. This is the potential relevant for both Newton's gravitational force and the Coulomb force of electrostatics. 


\para
{\co Recall that the centre of mass is
%
\be M{\bf R} = m_1\vecx_1+m_2\vecx_2\ .\ee
%
We  also define the relative separation
%
\be {\bf r} = \vecx_1-\vecx_2\ .\ee
%
With a potential $V(r)$, the equations of motion are
%
\be m_1\ddot{\bx}_1 = -\ppp{V}{\bx_1} = -\ppp{V}{r}\hat{\bf r}\ \ \ {\rm and}\ \ \ m_2\ddot{\bx}_2 = -\ppp{V}{\bx_2}=+\ppp{V}{r}\hat{\br}\ .\ee
%
This tells us that $\ddot{R}=0$} so that the centre of mass travels with constant velocity. This reaffirms a result we saw previously if there is no external force, then the total momentum of the system is conserved. For us, it  simply means that we can ignore the centre of mass coordinate. 

\para
{\co We also have
%
\be \ddot{\br} = \ddot{\bx}_1-\ddot{\bx}_2 = -\left(\frac{1}{m_1} +\frac{1}{m_2}\right)\ppp{V}{r}\hat{\br}\ .\label{tomdidthis}\ee
%
We write this as 
%
\be \mu\ddot{{\bf r}} = -\nabla V\label{station}\ee
%
where we have defined the {\it reduced mass}
%
\be \mu =\frac{m_1m_2}{m_1+m_2}\ .\ee
%
}
But this is exactly the kind of single-particle problem that we looked at when we started these lectures.  The only novelty is that the mass $m$ gets replaced with the reduced mass $\mu$. 

\para
In what follows, we'll sometimes employ the language of one-particle mechanics and talk about the ``the particle with mass $\mu$ and  position $\br$". That's really shorthand for the positions of two particles, with relative separation $\br$ and reduced mass $\mu$. 

\para
{\co In the limit where one of the particles involved is very heavy, say $m_2\gg m_1$, then the reduced mass is approximately equal to the lighter of the two: $\mu \approx m_1$. In this case, the heavy object is essentially fixed, with the lighter object orbiting around it.} This was the approximation that we were implicitly invoking earlier we referred to the lighter particle as a ``test particle". It's a very good approximation for many cases of interest. For example, the centre of mass of the Earth and Sun is very close to the centre of the Sun and it makes sense to think of the Sun as fixed, and the Earth moving in its gravitational field.  Even for the Earth and Moon, the centre of mass lies 1000 miles below the surface of the Earth.


\subsection{Learning Objectives}

Students should appreciate how we can extract the centre of mass, so that the two-body problem really becomes a one-body problem. 

\newpage
\section{Conservation of Angular Momentum}

{\co We have a potential of the form
%
\be V(|\bx_1-\bx_2|) = V(r)\ .\ee
%
But for such central potentials, the angular momentum
%
\be 
{\bf L} = \mu\vecr\times\dvecr\ \label{angmom}\ee
%
is conserved. This follows from  taking the time derivative
%
\be \frac{d{\bf L}}{dt} = \mu\vecr\times\ddvecr =  -\vecr\times \nabla V\ee
%
and noting that, for a central potential $V=V(r)$, the gradient $\nabla V$ is parallel to $\br$.}

\para
This implies that the motion takes place in a plane perpendicular to ${\bf L}$ . This follows because ${\bf L}$ is a fixed, unchanging vector which, by construction, obeys 
%
{\co \be {\bf L}\cdot \vecr=0\ .\ee
%
So the relative position of the two particles  always lies in a plane perpendicular to ${\bf L}$. } By the same argument,  ${\bf L}\cdot \dvecr=0$ so the relative velocity of the particles also lies in the same plane. In this way the three-dimensional dynamics is reduced to dynamics on a plane. 

\para
We've learned that the motion lies in a plane. It will turn out to be much easier if we work with polar coordinates in the plane rather than Cartesian coordinates.  For this reason, we take  a brief detour to explain some relevant aspects of polar coordinates. 



\para
To start, we rotate our coordinate system so that the angular momentum points in the $z$-direction and all motion takes place in the $(x,y)$-plane.  {\co Pick ${\bf L}$ to point in the $\hat{\bf z}$-direction. We then define the usual polar coordinates
%
\be x = r\cos\theta\ \ \ {\rm and}\ \  \ y = r\sin\theta\ . \ \ \ \ \ee
%
}
Our goal is to express both the velocity and acceleration in polar coordinates. {\co Itroduce two unit vectors, $\hat{{\bf r}}$ and $\hat{{\btheta}}$ in the direction of increasing $r$ and $\theta$ respectively, as shown in the {\cblue draw diagram}. Written in Cartesian form, these vectors are
%
\be \hat{{\bf r}} = \left(\begin{array}{c}\cos\theta \\ \sin\theta \end{array}\right)\ \ \ {\rm and}\ \ \  \hat{{\btheta}} = \left(\begin{array}{c}-\sin\theta \\ \cos\theta \end{array}\right)\ .\ee
%
}
Note that we've had a subtle change of notation, since $\hatr$ is now a 2d vector, rather than the 3d vector. This is appropriate because the motion lies in the a plane.

\para
The vectors $\hatr$ and $\hat{\btheta}$ form an orthornormal basis at every point in the plane. But the basis itself depends on the point $\br=(x,y)$ at which we sit. Or, more specifically, the basis depends on the angle $\theta$: moving in the radial direction doesn't change the basis, but moving in the angular  direction does. {\co We have
%
\be \frac{d\hat{{\bf r}}}{d\theta} = \left(\begin{array}{c}-\sin\theta \\ \cos\theta \end{array}\right) = \hat{{\btheta}}\ \ \ {\rm and}\ \ \ \frac{d\hat{{\btheta}}}{d\theta} = \left(\begin{array}{c}-\cos\theta \\ -\sin\theta \end{array}\right) = -\hat{{\bf r}}\ .\label{gethatson}\ee
%
}
This means that if the particle moves in a way such that $\theta$ changes with time, then the basis vectors themselves will also change with time. This has an implication for the velocity when expressed in polar coordinates.  To see this, we first {\co write the position of the particle as the simple, if somewhat ugly, equation
%
\be {\bf r} = r\hat{{\bf r}}\ .\ee
%
From this we can compute the velocity, remembering that both $r$ and the basis vector $\hat{{\bf r}}$ can change with time. We get
%
\be \dvecr =  \dot{r}\hatr+r\frac{d\hatr}{d\theta}\dot{\theta} =\dot{r}\hatr+r\dot{\theta}\hattheta\ .\label{polarv}\ee
%
}
The second term in the above expression arises because the basis vectors change with time, and is proportional to the... {\co Here  $\dot{\theta}$ is the {\it angular velocity}.}  (Sometimes this is referred to as the {\it angular speed},  even though it can be positive or negative, with the term ``angular velocity" reserved for a 3d vector $\bomega$ that we met previously. 

\para
{\co 
Differentiating once more gives us the expression for acceleration in polar coordinates
%
\be \ddvecr &=& \ddot{r}\hatr + \dot{r}\frac{d\hatr}{d\theta}\dot{\theta} +\dot{r}\dot{\theta}\hattheta + r\ddot{\theta}\hattheta+r\dot{\theta}\frac{d\hattheta}{d\theta}\dot{\theta} \nn\\ &=& (\ddot{r}-r\dot{\theta}^2)\hatr+(r\ddot{\theta}+2\dot{r}\dot{\theta})\hattheta\ .\label{polara}\ee
%
}
The two expressions \eqn{polarv} and \eqn{polara} will be important in what follows.

\subsubsection*{An Example: Circular Motion}

We can illustrate these expressions with something familiar: circular motion. {\co A particle moving in a circle has $\dot{r}=0$. If the particle travels with constant angular velocity $\dot{\theta}=\omega$ then, from \eqn{polarv}, the velocity in the plane is
%
\be \dvecr = r\omega \hattheta\ .\ee
%
The speed in the plane is then $v=|\dvecr| = r\omega$. Similarly, from \eqn{polara}, the acceleration in the plane is
%
\be \ddvecr = -r\omega^2\hatr\ .\ee
%
The magnitude of the acceleration is $a=|\ddvecr| = r\omega^2=v^2/r$. This means that,  if we want a particle to travel in a circle, then we need to supply a force $F=mv^2/r$ towards the origin. This is known as a {\it centripetal force}.}


\subsection*{Comments}

{\cblue This lecture is largely setting up the formalism. Perhaps we can ask some simple questions on circular motion? Not sure what else.}

\subsection{Learning Objectives}

Students should be comfortable with working in polar coordinates in the plane, and how the equations of motion look in these coordinates. 


\newpage
\section{The Effective Potential}

We can now return to our central force problem, armed with polar coordinates. We know that the particle motion takes place in a plane, and we choose to parameterise this by the polar coordinates $r$ and $\theta$. {\co In the last lecture, we derived an expression for the velocity
%
\be 
\dvecr =  \dot{r}\hatr+r\frac{d\hatr}{d\theta}\dot{\theta} =\dot{r}\hatr+r\dot{\theta}\hattheta\ee
%
and the  acceleration in polar coordinates
%
\be \ddvecr  &=& (\ddot{r}-r\dot{\theta}^2)\hatr+(r\ddot{\theta}+2\dot{r}\dot{\theta})\hattheta\ .\ee
%
where 
%
\be \hat{{\bf r}} = \left(\begin{array}{c}\cos\theta \\ \sin\theta \end{array}\right)\ \ \ {\rm and}\ \ \  \hat{{\btheta}} = \left(\begin{array}{c}-\sin\theta \\ \cos\theta \end{array}\right)\ .\ee
%
Newton's second law then becomes
%
%
\be \mu(\ddot{r}-r\dot{\theta}^2)\hatr+\mu(r\ddot{\theta}+2\dot{r}\dot{\theta})\hattheta =-\frac{dV}{dr}\hatr\ .\label{centralpolar}\ee
%
with $\mu$ the reduced mass.

\para
The $\hattheta$ component of this equation is particularly simple. It reads
%
\be r\ddot{\theta} +2\dot{r}\dot{\theta} =0\ \ \ \Longrightarrow\ \ \  \frac{1}{r}\frac{d}{dt}\left(r^2\dot{\theta}\right)=0\ .\ee
%
At first, it looks as if we've found a new conserved quantity, since this equation is telling us that
%
\be l = r^2\dot{\theta}\label{l}\ee
%
does not change with time.} However, we shouldn't get too excited. This is something that we already know. To see this, let's look again at the angular momentum ${\bf L}$. We already used the fact that the direction of ${{\bf L}}$ is conserved, restricting motion to the plane. But what about the magnitude of ${\bf L}$? {\co The angular momentum is 
%
\be {\bf L} = \mu\vecr\times\dvecr= \mu r\hatr \times \left(\dot{r}\hatr+r\dot{\theta}\hattheta\right) = \mu r^2\dot{\theta}\,\left(\hatr\times \hattheta\right)\ .\ee
%
Since $\hatr$ and $\hattheta$ are orthogonal unit vectors, $\hatr\times\hattheta$ is also a unit vector, pointing out of the plane. The magnitude of the angular momentum vector is therefore
%
\be |{\bf L}| = \mu l\ .\ee
%
}
This means that $l=r^2\dot{\theta}$, is identified as the angular momentum per unit mass, although we will often be lazy and refer to $l$ simply as the angular momentum. 

\para
{\co Next, we look at the $\hatr$ component of the equation of motion
%
\be \mu(\ddot{r}-r\dot{\theta}^2) = -\frac{dV}{dr}\ .\ee
%
Using  $l=r^2\dot{\theta}$ is conserved, we can write this as
%
\be 
\mu\ddot{r} = -\frac{dV}{dr} + \frac{\mu l^2}{r^3}\ .\label{pinky}\ee
%
}
It's worth pausing to reflect on what's happened here. We started with two particles, each moving in three dimensions. This means that the system was described by six coordinates, $\bx_1$ and $\bx_2$.  We used momentum conservation to reduce this to a problem of just three coordinates, the relative separation $\br=\bx_1-\bx_2$.   We then used the direction of the angular momentum to reduce it to a two-dimensional problem, and the magnitude of the angular momentum to reduce it to a one-dimensional problem \eqn{pinky}, involving only the variable $r$. 



\para
 This should give you some idea of how important conserved quantities are when it comes to solving anything.   Roughly speaking, this is also why it's not usually possible to solve problems with many particles when they are mutually interacting.   When we have $N\geq 3$ particles, each with a mutual interaction potential, we typically don't  get any more conserved quantities beyond those that we've seen for the two-body problem. Without additional constraints, the motion becomes much more complicated. Even the simplest, the so-called {\it three-body problem}, exhibits chaotic behaviour. 
 

\para
{\co We have
%
\be \mu \ddot{r} = -\frac{dV_{\rm eff}}{dr}\ \label{eomeff}\ee
%
where $V_{\rm eff}(r)$ is called the {\it effective potential}
%
\be V_{\rm eff}(r) = V(r) +\frac{\mu l^2}{2r^2} \ .\label{veff}\ee
%
The extra term, $\mu l^2/2r^2$ is called the {\it angular momentum barrier}}.  (also known as the centrifugal barrier). As you can see, this extra term diverges as  $r\ra 0$ and the two particles get close.



\para
To get more intuition for the effective potential, we can look at the total energy of the system.{\co The total energy is
%
\be E&=& \frac{1}{2}m_1\dvecx_1\cdot\dvecx_1 + \frac{1}{2}m_2\dvecx_2\cdot\dvecx_2 + V(|\bx_1-\bx_2|)\nn\\
&=& \frac{1}{2}(m_1+m_2)\dot{\bR}^2 + \frac{1}{2}\mu \dot{\br}^2 + V(r) \ .\ee
%
We'll ignore the kinetic energy for the centre of mass since that decouples from the interesting dynamics. 
If we evaluate the remaining energy  for some fixed angular momentum ${\bf L}$ then}  using the $\dot{\bf r}$ equation in polar coordinates and the fact that the particle moves only in the plane, {\co we have
%
\be E &=& \frac{1}{2}\mu\dot{r}^2 + \frac{1}{2}\mu r^2\dot{\theta}^2 + V(r) \nn\\ &=& \frac{1}{2}\mu\dot{r}^2 +\frac{\mu l^2}{2r^2}+V(r)\label{rainstopped}\\ &=& \frac{1}{2}\mu \dot{r}^2 + V_{\rm eff}(r)\ .\nn\ee
%
}
We now clearly see that the extra term in the effective potential, the angular momentum barrier, has its origin as the angular part of the kinetic energy. As the particles get close, the conservation of angular momentum means that their angular velocity must increase. And this has the effect of increasing the kinetic energy. We  include this ``angular kinetic energy" in the effective potential, where it prohibits the particles from getting too close, now by paying a heavy price in ``effective potential energy".

\subsection{Learning Objectives}

Students should be appreciate how an extra term in the potential can account for the angular motion. 

\newpage
\section{Getting a Feel for Orbits}


In writing our equation of motion in the form \eqn{eomeff}, we've succeeded in reducing our original 3d problem to a 1d problem in the radial coordinate $r$. But we've already seen how to solve 1d problems of this kind in the first part of these lectures. In particular, we saw that we can understand qualitative aspects of one-dimensional motion simply by plotting the potential energy and thinking about how the particle rolls about. Here we do the same thing for the effective potential.



\para
We can play this game for any central potential $V(r)$, but we will  focus our attention on the most important. {\co Consider
%
\be V(r)=-\frac{\mu k}{r}\ .\ee
%
}
The factor of $\mu$ in this expression is for convenience as it ensures that $\mu$ drops out of many subsequent equations. {\co For the mutual gravitational interaction, we should take $k=Gm_1m_2/\mu = G(m_1+m_2)$. Meanwhile, for the Coulomb force we should take $k=-q_1q_2/4\pi\epsilon_0\mu$.  In both cases, the resulting inverse square law is attractive for $k>0$ and repulsive for $k<0$.
 The effective potential is 
%
\be V_{\rm eff}(r) = -\frac{\mu k}{r}+\frac{\mu l^2}{2r^2}\ .\ee
%
{\cblue Draw figure for $k>0$.}} From this potential, we can easily learn about the motion $r(t)$ in the radial direction. The tricky part is to superpose this with the angular velocity, $\dot{\theta}= l/r^2$, which is fixed (as a function of $r$) by angular momentum conservation. In this way, we build up the full 3d picture of the particle's trajectory.





\para
The possible forms of the motion can be characterised by their energy $E$:
%

\begin{itemize}{\co
\item $E=E_{\rm min}=-\mu k^2/2l^2$: Here the particle sits at the minimum of the effective potential,
%
\be r_\star=\frac{l^2}{k} \ee
%
and stays there  for all time. But $\dot{\theta} = l/r_\star^2$ is constant so the particle moves in a circle.}



\vskip 3mm
Notice that the position $r_\star=l^2/k$ of the minimum depends on the angular momentum $l$. The higher the angular momentum, the further away the minimum. If there is no angular momentum, and $l=0$, then $V_{\rm eff}(r)= V(r)$ and the potential has no minimum. This is telling us the obvious fact that there is no way that $r$ can be constant unless the particle is moving in the $\theta$ direction.



{\co 
\item  $E_{\rm min}<E<0$: Here $r$ oscillating back and forth. Again, because $\dot{\theta}\neq 0$, the particle also moves in the angular direction. The result is a non-circular orbit.}  Although it is not yet obvious, we will soon show that for $V=-\mu k/r$, this orbit is an ellipse.

\vskip 3mm
{\co The point on the orbit with the smallest value of $r$ is called the {\it periapsis} (or {\it perihelion} for motion around the Sun). The point with the largest value is the {\it apoapsis} (or {\it aphelion} for the Sun.).} Together, these two points are referred to as the {\it apsides}.  

\vskip 3mm
{\co \item $E>0$. Now the particle can escape to $r\rightarrow \infty$. This is not an orbit.}  If you send it in with some velocity from infinity, then it navigates the dip, moves up towards the origin before reaching some minimum distance, and  then rolls back out to infinity. We will see later that, for the $V=-\mu k/r$ potential, the resulting trajectory is a hyperbola. 
\end{itemize}

\vskip 2mm
\subsection{Stability of Circular Orbits}

Before we learn how to solve for general orbits, we first pause to ask some general questions. Given some potential $V(r)$, we can ask: when do circular orbits exist? And when are they stable?

\para
{\co Circular orbits are equilibrium points such that
%
\be V^\prime_{\rm eff}(r_\star) =0\ .\label{centralcritical}\ee
%
The orbit is stable if small perturbations return us back to the critical point, so
%
\be V_{\rm eff}^{\prime\prime}(r_\star)>0\ .\ee
%
}
If this condition holds, small radial deviations from the circular orbit will oscillate about $r_\star$ with simple harmonic motion. 

\para
{\co For example
%
\be V(r) = -\frac{\mu k}{r^n}\  \ \ {\rm with}\ \  n\geq 1\ee
%
always has circular orbits for  $k>0$, sitting at a radius $r_\star^{2-n} = l^2/nk$. You can check
%
\be V^{\prime\prime}(r_\star) +  \frac{3}{r_\star}V^\prime(r_\star) = -\Big(n(n+1) - 3n\Big)\frac{\mu k}{r_\star^{n+2}} >0\ .\ee
%
This holds only for $n<2$. {\cblue plot figures for $n=1$ and $n=3$.}




\para
Here's a cute fact: in a universe with $d$ spatial dimensions, the law of gravity would be $F\sim 1/r^{d-1}$ corresponding to a potential energy $V\sim -1/r^{d-2}$. Stable circular orbits only exist in $d=\leq 3$ dimensions!}

\subsection*{Comments}

{\cblue It would be nice to try to get a problem built from this lecture and the previous one. One of the tricky things is trying to piece together the 1d motion in the $r$ direction, with the motion in the $\theta$ direction fixed by angular momentum. It would be good to have something that fleshes this out.}

\subsection{Learning Objectives}

Students start to understand how motion in the angular direction is tied in with motion in the radial direction. 




\newpage
\section{Orbits}

In this lecture, we're going to bite the bullet and solve for the motion of two objects -- say a star and a planet -- moving under the Newtonian $1/r^2$ force law. 


  
 \para
We know that the motion is in a plane, and we've already intimated that polar coordinates $r$ and $\theta$ are useful. The obvious thing to do is to   solve for $r(t)$ and $\theta(t)$. That turns out to be tricky. Instead, the right approach is to try to find an expression for the shape of the orbit....{\co We will solve for $r=r(\theta)$. 

\para
In fact, we will work with 
%
\be u = \frac{1}{r}\ .\ee
%
}
I wish I had a reason to motivate this trick. Unfortunately, I don't. You'll just have to trust me and we'll see that it helps. We will solve for $u(\theta)$. 

\para
{\co The velocity is
%
\be \frac{dr}{dt} = \frac{dr}{d\theta} \dot{\theta} = \frac{dr}{d\theta} \frac{l}{r^2} = -l\frac{du}{d\theta}\ \ee
%
with $l=r^2\dot{\theta}$ the conserved angular momentum. 
Meanwhile, the acceleration is 
%
\be \frac{d^2r}{dt^2} = \frac{d}{dt}\left(-l\frac{du}{d\theta}\right)  = -l\frac{d^2 u}{d\theta^2}\dot{\theta} = -l^2\frac{d^2u}{d\theta^2}\frac{1}{r^2} = -l^2u^2\frac{d^2u}{d\theta^2}\ .\label{uac}\ee
%
The equation of motion for the radial position is
%
\be \mu \ddot{r} - \frac{\mu l^2}{r^3} = F(r)\ee
%
with  $F(r) = -dV/dr$.}  Using \eqn{uac}, and doing a little bit of algebra (basically dividing by $\mu l^2u^2)$, {\co we get 
%
\be \frac{d^2u}{d\theta^2} + u = -\frac{1}{\mu l^2u^2} F(1/u)\ .\label{orbit}\ee
%
This is the {\it orbit equation}.} Our goal is to solve this for $u(\theta)$. If we want to subsequently figure out the time dependence, we can then extract it from the equation  $\dot{\theta} = lu^2$.

\subsection{The Kepler Problem}

The {\it Kepler problem}\index{Kepler problem} is the name given to understanding planetary orbits about a star. It is named after the astronomer Johannes Kepler who, through many years of astronomical observation, uncovered the mathematical patterns evident in our solar system. We will describe Kepler's observations shortly.

\para
{\co We have the potential 
%
\be V(r)  = -\frac{\mu k}{r}\label{burnt}\ee
%
The orbit equation is
%
\be \frac{d^2u}{d\theta^2} +u = \frac{k}{l^2}\ .\label{orbitfork}\ee
%
But this is just the equation for a harmonic oscillator, albeit with its centre displaced by $k/l^2$! We can write the most general solution as 
%
\be u = C\cos\left(\theta-\theta_0\right) + \frac{k}{l^2}\label{thankwells}\ee
%
with $C\geq 0$ and $\theta_0$ integration constants. } You might be tempted instead to write $u=C\cos\theta + C'\sin\theta +k/l^2$ with $C$ and $C'$ as integration constants. This is equivalent to our result above but, as we will now see, it's much more useful to use $\theta_0$ as the second integration constant.)

\para
{\co 
At the point where the orbit is closest to the origin (the periapsis), $u$ is largest. We  have $u_{\rm max} = C+k/l^2$. We choose to orient our polar coordinates so that the periapsis occurs at $\theta=0$. This choice means that we  set $\theta_0=0$. In terms of our original variable $r=1/u$, we then have the final expression for the orbit
%
\be r = \frac{r_0}{e\cos\theta + 1}\label{roundandround}\ee
%
where
%
\be r_0 = \frac{l^2}{k}\ \ \  {\rm and}\ \ \  e = \frac{Cl^2}{k}\ . \label{rzeroande}\ee
%
}
Notice that $r_0$ is fixed by the angular momentum, while the choice of $e$ is now effectively the integration constant in the problem. We necessarily have $e\geq 0$.

\para
It turns out that  this is a well-known equation in geometry: it describes a geometrical object called  a {\it conic section}\index{conic section}. These are the curves that you get if you intersect a cone with some plane. As we will now derive in some detail, the resulting curves -- which, for us, are orbits -- can be either ellipses (including circles as a special case), hyperbolae, or parabolae. Which curve we get is determined by the  integration constant $e$ which is called the {\it eccentricity}\index{eccentricity}. In particular, it matters a lot whether $e<1$ or $e\geq 1$. We look at these in turn.


{\co 
\subsubsection*{Ellipses: $e<1$}

For $0<e<1$, the radial position  is bounded in the interval
%
\be \frac{r_0}{r}\in[1-e,1+e]\ .\ee
%
We can convert \eqn{roundandround} back to Cartesian coordinates, $x=r\cos\theta$ and $y=r\sin\theta$, by writing
%
\be r=r_0-er\cos\theta\ \ \  \Longrightarrow   \ \ \ x^2+y^2 = (r_0-ex)^2\ .\ee
%
Multiplying out the square, collecting terms, and rearranging allows us to write this equation as the equation of an {\it ellipse}\index{ellipse}
%
\be \frac{(x-x_c)^2}{a^2}+\frac{y^2}{b^2}=1\ \ee
%
with
%
\be x_c=-\frac{er_0}{1-e^2}\ \ \ ,\ \ \ \ a^2 = \frac{r_0^2}{(1-e^2)^2}\ \  \ ,\ \ \ b^2 = \frac{r_0^2}{1-e^2} < a^2 \ .\label{aandb}\ee
%
%\EPSFIGURE{ellipse.eps,height=70pt}{}
%
{\cblue Draw figure} }The two semi-axes  have lengths $a$ and $b$, while its centre is shifted from the origin by a distance  $x=x_c<0$.   This shift is important. The origin at $r=0$ is the centre of attraction of the gravitational potential. 




\para
In general, the origin should be thought of as the centre of mass of the two orbiting particles. But, as we mentioned previously, the centre of mass of the Earth and Sun system is essentially the same as the centre of the Sun.  So it's appropriate to think of the origin as the location of the Sun. But this does not coincide with the centre of the elliptical orbit.  {\co The Sun is centred at the origin of our coordinate system. The distance to the centre of the ellipse is
%
\be |x_c|=\frac{r_0e}{1-e^2} = ea\ .\ee
%
The point where the Sun sits is called the {\it focus} of the ellipse.} In fact, it is one of two foci, one denoted $O$ and the other $O'$ in the right-hand figure, both distance $|x_c|$ from the centre along the major axis.  {\cblue maybe mention geometric meaning of two foci.}

\para
{\co When $e=0$, the focus sits at the centre of the ellipse and the lengths of the two axes coincide: $a=b$. This is a circular orbit.} 

\para
{\co In the Solar System, nearly all planets have $e<0.1$. This means that the difference between the major and minor axes of their orbits are equal to within less than $1\%$ of each other, and so the orbits are very nearly circular. The only exception is Mercury, the closest planet to the Sun, which has $e\approx 0.2$.  There are, however, objects that lie on very eccentric orbits. These are comets. The most famous is Halley's comet, with  eccentricity $e\approx 0.97$.




\subsubsection*{Hyperbolae: $e>1$}



For $e>1$, the orbit is no longer bound. There  are two values of $\theta$ for which $r\rightarrow \infty$. This occurs when  $\cos\theta = -1/e$.} Repeating the algebraic steps that lead to the equation of an ellipse, {\co we now find
%
\be \frac{(x-ae)^2}{a^2}-\frac{y^2}{b^2}=1\ \label{hyperbola}\ee
%
 with 
 %
 \be a^2 = \frac{r_0^2}{(e^2-1)^2} \ \ \ {\rm and} \ \ \ b^2 = \frac{r_0^2}{e^2-1}\ .\ee
 %
 This is the equation for a {\it hyperbola}.}
That relative minus signs between the two terms in \eqn{hyperbola} is all important. {\cblue draw figure}. It is plotted in the figure, where the dashed lines are the asymptotes. They meet at the point $x=r_0e/(e^2-1)$. Again, the centre of the gravitational attraction sits at the origin denoted by the sun-like disc. 

\para
We've already noted that  the orbit goes off to $r\rightarrow\infty$ when $\cos\theta = -1/e$. Because $-1/e$ is negative, this must occur for two angles in the range $\pi/2< \theta < 3\pi/2$. This is one way to see why the orbit sits in the left-hand quadrant as shown in the figure.

{\co 
\subsubsection*{Parabolae: $e=1$}

Finally, in the special case of $e=1$, the algebra is particularly simple. The orbit is described by 
%
\be y^2 = r_0^2 - 2r_0 x\ .\ee
%
This is the equation for a parabola.}

\subsubsection*{A Repulsive Force}

{\co Above, we assumed that $k>0$, meaning the force is attractive. For $k<0$, we get
%
%
\be r = \frac{r_0}{e\cos\theta - 1}\ .\ee
%
where $r_0 = l^2/|k|$ and $e=Cl^2/|k|$ with $C\geq 0$. } Note that with this choice of convention, we retain  the condition $e\geq 0$. But, because we must necessarily have $r>0$, we only find solutions in  the case $e>1$. This is nice: we wouldn't expect to find bound orbits between two particles which repel each other. For $e>1$, the unbounded orbits are again hyperbolae and look like that shown in {\cblue draw figure}. Notice that the orbits go off to $r\rightarrow \infty$ when $\cos\theta = +1/e$ which, since $e>1$, must occur at an angle $0<\theta <\pi/2$ and an angle $3\pi/2<\theta<2\pi$.  This is the reason that the orbit sits in the right-hand quadrant. 


\subsection*{Comments}

{\cblue 
There's an obvious opportunity for simulations or orbits. Perhaps another one demonstrating how these orbits are related to slicing the cone. Maybe others? What about problems? Surely there are standard things out there?}

\para
{\cblue One possibility for a problem, albeit one that is slightly advanced, is to show that the Runge-Lenz vector is conserved. This gives the conic section equation immediately. I've included the calculation below. 


\subsection{Learning Objectives}

Students should be able to derive the possible motions of planets and comets from first principles. 


\subsection{A Possible Problem: The Runge--Lenz Vector}

There is another, more elegant, way to see that the orbit of a planet traces out an ellipse. This is because a particle of mass $\mu$ moving under the inverse-square force
%
\be {\bf F} = -\frac{\mu k}{r^2}\hat{\br}\ \ee
%
enjoys another conserved quantity, namely
%
\be {\bf A} = \frac{\dot{\bx}\times {\bf L}}{\mu k} - \hat{\br}\ \label{rungelenz}\ee
%
where $\hat{\br} = \bx/r$ and we've reverted to notation that we previously used for a test particle, where the Sun sits at the origin and the planet has position $\bx$.  The vector ${\bf A}$ is known as the {\it Runge--Lenz vector}, or sometimes as the {\it Laplace--Runge--Lenz vector}. (In fact, with this normalisation ${\bf A}$ is more properly called the {\it eccentricity vector}.

\para
The existence of an additional conserved quantity in the Kepler problem is extremely surprising. It is also very uncommon. Pretty much any other potential does not come with extra conservation laws. But the $V\sim 1/r$ potential is mathematically special. It is a wonderful fact, and one that is repeated throughout physics, that the formulae that have special mathematical properties are also those relevant for nature.

\para
To see that the Runge--Lenz vector is conserved, we simply need to differentiate. We use the fact that angular momentum is conserved, so $\dot{\bf L}=0$, to write
%
\be \dot{\bf A} = \frac{\ddot{\bx} \times {\bf L}}{\mu k}   -\frac{\dot{\bx}}{r} + \frac{\dot{r}}{r^2}\bx\ .\ee
%
Now $r^2= \bx\cdot\bx$, from which we learn that $\dot{r} = (\bx\cdot\dot{\bx})/r$. The above formula then becomes
%
\be \dot{\bf A} = \frac{\ddot{\bx} \times (\bx\times \dot{\bx})}{k} - \frac{\dot{\bx}}{r} + \frac{\bx\cdot\dot{\bx}}{r^3}\bx =  -\frac{\hat{\br} \times (\bx\times \dot{\bx})}{r^2} -   \frac{\dot{\bx}}{r} + \frac{\bx\cdot\dot{\bx}}{r^3}\bx\ .\ \ \ \ee
%
Now it's just a matter of using the triple product formula, ${\bf a}\times ({\bf b}\times{\bf c}) = ({\bf a}\cdot {\bf c}){\bf b} - ({\bf a}\cdot {\bf b}){\bf c}$ to see that the first term cancels the second two. This means that $\dot{\bf A}=0$. 

\para
The conservation of the Runge--Lenz vector has a number of consequences. But for now we can use it to quickly re-derive our earlier result. The vector ${\bf A}$ points in some fixed direction in space. Note that ${\bf A}\cdot {\bf L}=0$, which means that ${\bf A}$ points in some direction in the plane of the orbit. If we take the inner product with the position vector, we get
%
\be {\bf A}\cdot\bx= Ar\cos\theta\ .\label{rltheta}\ee
%
For now, we'll think of this equation as {\it defining} the angle $\theta$ in the plane. (We'll shortly see that it coincides with our previous definition.) Using \eqn{rungelenz}, we have
%
\be {\bf A}\cdot \bx  = \frac{\bx\cdot(\dot{\bx }\times {\bf L})}{\mu k} - r =\frac{{\bf L}\cdot(\bx\times\dot{\bx})}{\mu k} -r = \frac{l^2}{k}-r\ee
%
where, in the second equality, we've used the symmetry of the scalar triple product. Putting this result together with our definition \eqn{rltheta}, we have
%
\be Ar\cos\theta = \frac{l^2}{k}  - r\ \ \ \Longrightarrow\ \ \ r= \frac{r_0}{A\cos\theta +1}\ \label{rlconic}\ee
%
where $r_0=l^2/k$. This is precisely the equation for a conic section \eqn{roundandround} that we previously derived by solving the orbit equation.

\para
There are a few lessons to take away from this. First, comparing \eqn{rlconic} with \eqn{roundandround}, we see that the magnitude of the Runge--Lenz vector coincides with the eccentricity, $|{\bf A}| = e$. (This is the reason that ${\bf A}$ is sometimes referred to as the eccentricity vector.) Second, we see that the angle $\theta$ is the same in the two formulae. Previously we defined $\theta=0$ to be the direction of the perihelion, the closest point of the orbit. Here we defined $\theta=0$ to be the direction in which ${\bf A}$ points. Combining these, we get a geometrical intuition for the Runge--Lenz vector: it points from the Sun towards the perihelion, and its magnitude is equal to the eccentricity of the orbit.

\para
The existence of the conserved Runge--Lenz vector explains a fact that is so familiar we didn't even comment on it above: the orbits in the Kepler problem are closed. This means that each time we go around the Sun, the closest point always sits at the same angle $\theta$. This is a special property of the inverse-square law that doesn't hold in other potentials. 
}


\newpage
\section{The Energy of the Orbit Revisited}

We can tally our different  solutions  -- the ellipse, hyperbola, and paraboa, with the general picture of orbits that we built by looking at the effective potential. {\co The energy  of a given orbit is
%
\be E&=& \frac{1}{2}\mu\dot{r}^2 + \frac{\mu l^2}{2r^2} - \frac{k\mu}{r} \nn\\ &=& \frac{1}{2}\mu \left(\frac{dr}{d\theta}\right)^2\dot{\theta}^2 +\frac{\mu l^2}{2r^2}-\frac{k\mu}{r} \nn\\ &=& \frac{1}{2}\mu\left(\frac{dr}{d\theta}\right)^2\frac{l^2}{r^4}+\frac{\mu l^2}{2r^2}-\frac{k\mu}{r}
\ .\ee
%
The orbit obeys
%

\be r = \frac{r_0}{e\cos\theta + 1}\label{roundandround}\ee
%
So
%
\be\frac{dr}{d\theta} = \frac{r_0e\sin\theta}{(1+e\cos\theta)^2}\ .\ee
%
}
Substituting in, after a couple of lines of algebra we find that all the $\theta$ dependence vanishes in the energy (as it must, since the energy is a constant of the motion). {\co We are left with
%
\be E = \frac{\mu k^2}{2l^2}(e^2-1)\ .\label{orbitenergy}\ee
%
We can look at what this means for the different types of orbits:}
%
\begin{itemize}
{\co \item $e=0 \ \Longrightarrow\ E=-\mu k^2/2l^2$. } This coincides with the minimum of the effective potential $V_{\rm eff}$  which, as we previously understood, corresponds to a circular orbit.\vskip 3mm
{\co \item $0<e<1 \ \Longrightarrow\ E<0$: These are the trapped, or bound, orbits} that we now know are ellipses. 

{\cblue Will probably skip this...maybe it could be a problem. If we substitute our expression for the semi-major axis $a$ given in \eqn{aandb}, together with $l^2 = r_0k$ from \eqn{rzeroande}, we can rewrite the energy as
%
\be E = -\frac{\mu k}{2a}\ .\ee
%
This is the archaically named {\it vis-viva equation}, \index{vis-viva equation} latin for ``living force". It tells us that the energy of an elliptic orbit depends only on the semi-major axis, and not on the eccentricity. } \vskip 3mm
{\co \item $e=1\ \Longrightarrow\ E=0$: This is the special case of a parabola. \vskip 3mm
\item $e>1\ \Longrightarrow \ E>0$: These are the unbounded orbits that we now know are hyperbolae. }
\end{itemize}




\subsection{Learning Objectives}

How energies can be positive or negative (or zero) and how these correlate with the different shapes of orbits. 

%\noindent



\subsection{A Possible Problem: Scattering}

{\cblue We decided not to include this in the lectures. But it could, perhaps, be developed as a long problem.

\para
We start by setting the scene, and introduce  some ideas that hold for any scattering problem.  We restrict our attention to central potentials $V(r)$, and assume that $V(r)\rightarrow 0$ as $r\rightarrow \infty$. 




\para
We do our experiment and throw  the particle from a large distance which we will take to be $r\rightarrow \infty$. We want to throw the particle towards the origin, but our aim is not always spot on. The particle has some energy $E>0$. It comes in from infinity and will typically be deflected by some angle, and ultimately escape out to infinity again.  If the interaction is repulsive, we expect that its trajectory will look something like that {\cblue draw figure}. However, much of what we're about to say will hold whether the force is attractive or repulsive. 




\para
First, by energy conservation,  the initial speed = final speed. This is true because of our assumption that $V(r)\rightarrow 0$ as  $r\rightarrow \infty$, so there is no contribution from the potential energy. We call this initial and final speed $v$.

\para
For a central potential, angular momentum, ${\bf L} = m\vecx\times\dvecx$ is conserved. We can get an expression for $l=|{\bf L}|/m$ as follows: draw a straight line tangent to the initial velocity. The closest this line gets to the origin 
is  a distance $b$, known as the {\it impact parameter}. {\cblue draw diagram}. The modulus of the angular momentum is then 
%
\be l = bv\ .\label{lbv}\ee
%
If this equation isn't immediately obvious mathematically, the following cute way of thinking may convince you. Suppose that there was no force acting on the particle at all. In this case, the particle would indeed follow the horizontal dotted straight line shown in the diagram. When the particle is closest to the origin, its velocity $\dot{{\bf x}}$  is perpendicular to its position ${\bf x}$ and  its angular momentum is obviously $l=bv$.  But angular momentum is conserved for a free particle, so this must also be its initial angular momentum. 



\para
Now we go back to the situation with the central potential $V(r)$. The angular momentum remains constant throughout the motion and, by the same kind of argument, at the end is given by $l=b^\prime v$ where $b^\prime$ is the shortest distance from the origin to the exit asymptote as shown in diagram But since the angular momentum is conserved, we must have $b = b^\prime$: the incoming impact parameter coincides with the outgoing impact parameter. 



\subsection*{Rutherford Scattering}

Historically, the most important scattering experiment involved throwing a charged particle at the nucleus to uncover the structure of the atom. {\cblue say more about this experiment here...} This process is now known as {\it Rutherford scattering}.


\para
Take a particle of charge $q$ and mass $m$ and throw it at a fixed particle of charge $Q$. The relevant force is the Coulomb potential
%
\be V = \frac{qQ}{4\pi\epsilon_0 r}\ .\ee
%
This is a repulsive force if $q$ and $Q$ have the same sign. Happily, the Coulomb force is mathematically identical to the gravitational force, which means that we can just take our earlier equations and replace 
%
\be k=-\frac{qQ}{4\pi\epsilon_0m}\ .\ee
%
{\cblue draw diagram} For a given scattering trajectory, the particle is deflected through a total angle $\psi$. We want to understand how this angle $\psi$ depends on various initial conditions of the trajectory.  Also introduce the  angle $\alpha$ given by 
%
\be \psi = \pi - 2\alpha\ .\label{phialpha}\ee
%
All the hard work has already been done.   For a repulsive $V\sim 1/r$ potential, we have
%
\be r = \frac{r_0}{e\cos\theta - 1}\ .\ee
%
The  particle asymptotes to $r\rightarrow \infty$ when $\cos\theta=1/e$. Comparing to the figure, this gives
%
\be \cos\alpha = \frac{1}{e}\ . \label{scatalpha}\ee
%
For a repulsive potential, we always have $e\geq 1$. We take $\alpha$ in the range $\alpha \in [0,\pi/2)$. 

\para
To proceed, we look at our various expressions for the  energy. We throw the particle from $r\rightarrow\infty$ at speed $v$. This means that when the particle started its journey, it had $E=\ft12 m v^2$.   But, in the last lecture, we derived the expression
%
\be E = \frac{1}{2}mv^2 = \frac{mk^2}{2l^2} (e^2-1) = \frac{mk^2}{2l^2}\tan^2\alpha\ . \ee
%
The angular momentum in this expression can be replaced by $l=bv$, and we get
%
\be \tan\left(\frac{\psi}{2}\right) = \frac{|k|}{bv^2}\ .\label{nottissue}\ee
%
This tells us the scattering angle $\psi$ for an initial impact parameter $b$ and velocity $v$. Note that $\psi\rightarrow \pi$ as $b\rightarrow 0$. That's to be expected: if you score a direct hit on the Coulomb target then you will bounce back the way you came. Conversely, $\psi\rightarrow 0$ as $b\rightarrow \infty$. That's telling you that if you miss the target by a lot, you'll just sail on past as if nothing happened. }




\newpage
\section{Kepler's Laws of Planetary Motion}


In 1605, Kepler published three laws which are obeyed by all planets in the Solar System. These laws were the culmination of decades of careful, painstaking observations of the night sky, firstly by Tycho Brahe and later by Kepler himself. {\co Kepler's laws of motion are
\begin{itemize}
\item {\bf K1:} Each planet moves in an ellipse, with the Sun at one focus. \vskip 3mm
\item {\bf K2:} The line between the planet and the Sun sweeps out equal areas in equal times. \vskip 3mm
\item {\bf K3:} The period $T$ of the orbit is proportional to the radius${}^{3/2}$. Or, said more precisely, the period of the orbit is proportional to $a^{3/2}$, where $a$ is the semi-major axis of the ellipse. In particular, the ratio $T^2/a^3$ is independent of the eccentricity of the ellipse. 
\end{itemize}
%
}
Before we go on, a cute historical fact. Kepler himself introduced the term ``focus" for the geometrical point of an ellipse. He wanted a term to demonstrate where the Sun sits. And, in latin, ``focus" means ``fireplace". 

\para
These laws were the clues that Newton needed to formulate the inverse square law of gravity. {\cblue Cute story here about Hooke, Halley, Wren and Newton....could be a tangent maybe?} Now that we understand the mathematics behind orbits, we will see how Kepler's  laws can be derived from Newton's inverse-square law of gravity. 



\para
{\co  Kepler's second law is just conservation of angular momentum. It holds for any central potential. It follows straightforwardly by looking at {\cblue draw figure} In time $\delta t$, the area swept out is 
%
\be \delta A = \frac{1}{2}r^2\delta \theta \ \ \ \Longrightarrow \ \ \ \frac{dA}{dt}=\frac{1}{2}r^2\dot{\theta} = \frac{l}{2}\ \ee
%
which we know is constant. }


\para
{\co Kepler's third law folllows from $F\sim 1/r^2$ and some  dimensionsal analysis. For gravity the mass $m$ of the planet drops out. We have
%
\be [GM] = L^3T^{-2}\ .\ee
%
}  ($G$ and $M$ do not appear individually). {\co Using this to relate the radius $R$ of an orbit with its period $T$, we have
%
\be be T^2 \sim \frac{R^3}{GM}\ .\ee
%
} Said differently, Kepler's third law is really telling us that the gravitational force goes as $F\sim 1/r^2$/

\para
{\cblue The actual story is more complicated and I'll probably sweep it under the rug. An orbit is also characterised by its angular momentum $l$ with
%
\be [l] = L^2 T^{-1} \ .\ee
%
So, in general, we can have
%
\be T^2 = \frac{R^3}{GM}\,f\left(\frac{GMR}{l^2}\right)\ee
%
with $f$ some arbitrary function. For a circular orbit of radius $R$, we have $GMR/l^2=1$, so this function $f$ is just a constant.

\para
 For an ellipse,  the parameter $GMR/l^2$ is related to the eccentricity of the orbit. The more precise claim of Kepler's third law is that  if we take our measure of the orbit to be $R=a$, the semi-major axis, then the function $f$ is just a constant. To see this, note that the area of an ellipse is 
%
\be A = \pi ab = \pi a^2 \sqrt{1-e^2} = \frac{\pi r_0^2}{(1-e^2)^{3/2}}\ .\ee
%
Since area is swept out at a constant rate $l/2$, the time period is 
%
\be T = \frac{2A}{l} = \frac{2\pi r_0^2}{l(1-e^2)^{3/2}} = \frac{2\pi}{\sqrt{GM}}\left(\frac{r_0}{1-e^2}\right)^{3/2} = \frac{2\pi}{\sqrt{GM}}a^{3/2}\ee
%
This is the precise form of Kepler's third law.}



%\para
%To fix the unknown function $f$, we need the more precise form of Kepler's third law.  This tells us that, for an elliptical orbit, it's best to take $R=a$, the semi-major axis of the ellipse. The improved version of K3 then says that $T^2/a^3$ is independent of the eccentricity, and so $f$ must be a constant. We then have
%
%\be T^2 \sim \frac{a^3}{GM}\ .\ee
%
%We already saw a version of this in  . For a general elliptical orbit, we can  be more precise. %
%where we've used the expression \eqn{andb} to write the period in terms of the length $a$ of the  semi-major axis. The quantity in brackets indeed has dimensions of a length. But what length is it? In fact, it has a nice interpretation. Recall that the periapsis of the orbit occurs at $r_{\rm min} = r_0/(1+e)$ and the apoapsis at $r_{\rm max} = r_0/(1-e)$. It is then natural to  define the average radius of the orbit to be $R = \ft12(r_{\rm min}+r_{\rm max}) = r_0/(1-e^2)$. We have
%
%\be T = \frac{2\pi}{\sqrt{GM}}R^{3/2}\ .\ee
%


\para
We see that Kepler's second  law and third  law follow on rather general grounds. It is Kepler's first law, which is seemingly the most simple, that is the most difficult to show. That, of course, is what we spent the previous lecture deriving. 


\subsection{Learning Objectives}

Students should appreciate how Kepler's historic laws can now be understood as a consequence of Newton's law of gravity. 

\newpage
\section{Optional: Perihelion Precession in General Relativity}

As we've seen,  a planet orbiting the Sun traces out an ellipse. But this analysis neglects two other important facts. The first is that there are other planets in the Solar System, and these too will exert a gravitational pull. The second is that Newton's theory of gravity isn't actually correct: it is superseded by Einstein's theory of general relativity, which describes gravity in terms of the bending of space and time.



\para
Both of these additional effects change the orbit. But, it turns out, they don't change it much. The orbit of a planet  in our Solar System remains approximately an ellipse, but the ellipse {\it precesses}. This means that the orbit doesn't close, and the point at which the planet is closest to the Sun -- the perihelion -- sits at a slightly different angle each time, as shown in the figure. 

\para
For nearly all the planets, it turns out that the observed precession can be correctly accounted for within Newtonian gravity by the gravitational pull of the other planets, with Jupiter giving the largest correction. The exception to this statement is Mercury. This is the planet closest to the Sun, and so feels the strongest gravitational field. 

\para
The full theory of general relativity is rather complicated and will be covered in a separate volume (where we will also make an attempt at describing the corrections to perihelion precession due to other planets). For our purposes, however, we can make do with something much simpler. It turns out that much of  the  effect of the curvature of spacetime can be captured in a simple correction to the effective potential.

\para
{\co The effects of general relativity can be approximated by
%
\be V_{\rm eff}(r) = -\frac{km}{r} + \frac{ml^2}{2r^2} - \frac{kml^2}{c^2r^3}\ .\ee
%
with  $k=GM$ and $l=r^2\dot{\theta}$ (treating the planet as a test particle.} The first two terms coincide with those of the Kepler problem. The third term is new and, indeed, introduces something novel: the speed of light $c$. This extra term becomes unimportant  when the planet orbits at distances $r\gg GM/c^2$, which we recognise as roughly the Schwarzschild radius. This means that we might expect the effects of the additional term to become noticeable for the planet closest to the Sun. 



\para
{\co If we repeat the derivation of the orbit equation, using $u=1/r$, we get
%
\be \frac{d^2u}{d\theta^2} + u = \frac{k}{l^2}+ \frac{3ku^2}{c^2}\ .\label{orbitforku}\ee
%
}
We'll treat the extra term, which is again tagged by that factor of $c^2$, as a perturbation to the original orbit. {\co We write $u(r) = u_0+ u_1(r)$ where $u_0$ is  a circular Newtonian orbit with
%
\be  u_0 = \frac{k}{l^2}\ee
%
and $u_1 \ll u_0$. Then, expanding to leading order in $u_1$, we get 
%
\be &&\frac{d^2u_1}{d\theta^2} + u_1 = \frac{3ku_0^2}{c^2} + \frac{6ku_0u_1}{c^2} \nn\\ \Longrightarrow\ \ \ && \frac{d^2u_1}{d\theta^2} + \left(1-\frac{6k^2}{c^2l^2}\right)u_1 = \frac{3k^3}{c^2l^4}\ .\ee
%
}
The solution to this equation is very similar to that of the Kepler problem. Including the $u_0=k/l^2$ term, {\co it is
%
\be u(\theta) = A\cos\left(\sqrt{1-\frac{6k^2}{c^2l^2}}\,\theta\right) + \frac{k}{l^2} +\frac{3k^3}{c^2l^4}\ \ee
%
}  where we have once again chosen our polar coordinates so that the integration constant is $\theta_0=0$.  This equation  describes an approximate ellipse. {\co This is an ellipse that precesses.  To see this, note that the periapsis occurs whenever the $\cos$ term is 1. This first happens at $\theta=0$. But the next time round, it happens at
%
\be \theta = 2\pi\left(1-\frac{6k^2}{c^2l^2}\right)^{-1/2}\approx\ 2\pi\left(1 +\frac{ 3k^2}{c^2l^2}\right)\ .\label{periprecess}\ee
%
This means that the perihelion advances by an angle of $6\pi G^2M^2/c^2l^2$ each turn. This is precisely what's needed to account for the extra  precession of Mercury. }



\subsection*{Comments}

{\cblue Not entirely sure that we want to end on this lecture...it's a little fiddly and is really a nod for things to come in General Relativity. Maybe we finish on Kepler's laws. }








\end{document}
